\documentclass[11pt]{article}
\usepackage[margin=1in]{geometry}
\usepackage[T1]{fontenc}
\usepackage{amsmath,amssymb,amsthm}
\usepackage{graphicx}
\usepackage[hidelinks]{hyperref}

\newtheorem{theorem}{Theorem}
\newtheorem{lemma}{Lemma}

\newcommand{\eps}{\varepsilon}
\newcommand{\ot}{\widehat{\otimes}_{\varepsilon}}
\newcommand{\BM}{B_{M(K)}}
\newcommand{\SF}{S_{\mathbb F}}
\newcommand{\norm}[1]{\left\|#1\right\|}
\newcommand{\Ahom}{A_{\mathrm{hom}}}

\title{The \(C(K)\)-Predual Subcase of Question 6.7 in arXiv:2604.18157}
\author{fa\_banach\_001, agent\_lane\_15}
\date{2026-06-29}

\begin{document}
\maketitle

\section*{Source and Question}

The source paper is Ondrej--Spurny, \emph{Operators on injective tensor
products of \(L_1\)-preduals}, arXiv:2604.18157v1 \cite{ondrej-spurny}.  Let
\(X\) be an \(L_1\)-predual and \(E,F\) Banach spaces.  If
\[
  T:X\ot E\to F
\]
is strongly bounded, the paper constructs a canonical extension
\[
  S:C(B_{X^*},E)\to F,\qquad Sf=\int_{B_{X^*}} f\,dG,
\]
where \(G\) is the representing measure of \(T\).

Theorem 6.5(b) proves that, for pseudo weakly compact and Right completely
continuous operators, \(S\) having the property implies \(T\) has it; the
converse is proved when \(X\) is separable.  The paper asks whether the
separability assumption can be removed.

\begin{center}
\includegraphics[width=\textwidth]{figures/open_problem_crop.png}
\end{center}

\section*{Result}

\begin{theorem}
\label{thm:main}
Let \(K\) be a compact Hausdorff space and let \(X=C(K)\), over the real or
complex scalar field \(\mathbb F\).  Let \(E,F\) be Banach spaces and let
\[
  T:C(K)\ot E \simeq C(K,E)\to F
\]
be a strongly bounded operator.  Let
\[
  S:C(B_{X^*},E)=C(B_{M(K)},E)\to F
\]
be the canonical extension from Ondrej--Spurny.

Then, for each of the following two operator classes,
\[
  \hbox{pseudo weakly compact}, \qquad
  \hbox{Right completely continuous},
\]
\(T\) belongs to the class if and only if \(S\) belongs to the class.
\end{theorem}

Thus Question 6.7 has an affirmative answer for all \(C(K)\)-preduals, with no
separability or metrizability assumption on \(K\).

\section*{Proof Intuition}

For a general \(L_1\)-predual, the source paper's boundary projection \(D\)
need only send continuous functions to Borel homogeneous affine functions.
Separability is used to replace this Borel object by continuous affine
selections.  For \(X=C(K)\), however, \(D\) has a concrete continuous formula:
first average a function over the homogeneous extreme boundary
\(\{\alpha\delta_t:\alpha\in S_{\mathbb F},\,t\in K\}\), then integrate the
result against \(\mu\in B_{M(K)}\).  This gives a norm-one projection from
\(C(B_{M(K)},E)\) onto the copy of \(C(K,E)\).  The canonical extension
therefore factors as \(S=T P\).  Since bounded maps preserve Right-null and
Right-Cauchy sequences, the desired preservation properties are immediate.

\section*{The Projection}

Let \(m\) denote normalized Haar measure on \(\SF\) (counting measure on
\(\{-1,1\}\) in the real case).  For
\(f\in C(\BM,E)\), define
\[
  (Rf)(t)=\int_{\SF}\alpha^{-1}f(\alpha\delta_t)\,dm(\alpha),
  \qquad t\in K,
\]
and then
\[
  (Pf)(\mu)=\int_K (Rf)(t)\,d\mu(t),
  \qquad \mu\in \BM.
\]
Here the second integral is the usual vector integral of a continuous
\(E\)-valued function against the scalar regular measure \(\mu\).

\begin{lemma}
\label{lem:projection}
\(P\) is a norm-one projection from \(C(\BM,E)\) onto
\(\Ahom(\BM,E)\), which is the canonical copy of \(C(K)\ot E\simeq C(K,E)\).
\end{lemma}

\begin{proof}
The map \(R:C(\BM,E)\to C(K,E)\) is bounded with \(\norm R\le 1\), since it is
an average of restrictions to the compact set
\(\{\alpha\delta_t:\alpha\in\SF,t\in K\}\).  If \(g\in C(K,E)\), its affine
homogeneous extension to \(\BM\) is
\[
  \widehat g(\mu)=\int_K g\,d\mu .
\]
Thus \(P f=\widehat{Rf}\) belongs to \(\Ahom(\BM,E)\) and \(\norm P\le 1\).
If \(f=\widehat g\) is already affine homogeneous, then
\[
  f(\alpha\delta_t)=\alpha g(t),
\]
so \(Rf=g\) and \(Pf=\widehat g=f\).  Hence \(P\) is a projection onto
\(\Ahom(\BM,E)\).  Its norm is one unless the space is trivial, because it is
the identity on constants of the form \(\widehat g\) with \(\norm g=1\).
\end{proof}

\begin{lemma}
\label{lem:factor}
For the canonical extension \(S\) of \(T\), one has \(S=T P\).
\end{lemma}

\begin{proof}
It is enough to check after applying an arbitrary \(y^*\in F^*\).  Let
\(G_{y^*}\) be the scalarized \(E^*\)-valued representing measure of \(T\).
In the source paper, scalar components of \(G_{y^*}\) belong to the
antihomogeneous boundary measure space associated with \(X\).

For \(X=C(K)\), the source paper's projection \(D\) on scalar measures has the
following explicit function-side adjoint: for scalar
\(\varphi\in C(\BM)\),
\[
  (D\varphi)(\mu)=
  \int_K\left(\int_{\SF}\alpha^{-1}\varphi(\alpha\delta_t)\,dm(\alpha)\right)
  d\mu(t).
\]
This is exactly the scalar version of \(P\).  Since \(G_{y^*}\) is already
antihomogeneous, the source identity \(D\nu=\nu\) for antihomogeneous boundary
measures gives
\[
  \int_{\BM} \varphi\,d\nu=\int_{\BM} D\varphi\,d\nu .
\]
Applying this to all scalar functions \(e^*\circ f\), \(e^*\in E^*\), yields
\[
  y^*(Sf)=\int_{\BM} f\,dG_{y^*}
          =\int_{\BM} Pf\,dG_{y^*}
          =y^*(T(Pf)).
\]
Since \(y^*\) was arbitrary, \(Sf=T(Pf)\).
\end{proof}

\section*{Proof of Theorem \ref{thm:main}}

\begin{proof}
The implication \(S\Rightarrow T\) is the easy half already present in
Theorem 6.5(b) of \cite{ondrej-spurny}, since \(T\) is the restriction of
\(S\) to \(\Ahom(\BM,E)\).

Conversely assume first that \(T\) is pseudo weakly compact.  If
\((f_n)\) is Right-null in \(C(\BM,E)\), then \((P f_n)\) is Right-null in
\(\Ahom(\BM,E)\), because \(P\) is bounded and bounded operators pull weakly
compact subsets of the target dual back to weakly compact subsets of the
domain dual.  Hence
\[
  \norm{S f_n}=\norm{T(Pf_n)}\to 0.
\]
Thus \(S\) is pseudo weakly compact.

The Right completely continuous case is identical.  If \((f_n)\) is
Right-Cauchy in \(C(\BM,E)\), then \((P f_n)\) is Right-Cauchy in
\(\Ahom(\BM,E)\).  Since \(T\) is Right completely continuous,
\((T(P f_n))\) converges in the Right topology of \(F\).  By Lemma
\ref{lem:factor}, this is exactly \((S f_n)\).  Therefore \(S\) is Right
completely continuous.
\end{proof}

\section*{Verification and Limitations}

The proof was checked against the source TeX around Lemma 2.7, Theorem 6.5,
and Question 6.7.  Local run indexes were searched for arXiv:2604.18157,
the paper title, Question 6.7, pseudo weakly compact, Right completely
continuous, \(C(K)\), and injective tensor product terms.  The run already
contains one partial packet for metrizable boundary support and one
no-promotion full-push attempt; no packet for this \(C(K)\)-predual
factorization subcase was found.

This is a genuine subcase only.  It does not settle arbitrary \(L_1\)-preduals,
because for a general \(L_1\)-predual the projection \(D\) need not be
continuous on \(C(B_{X^*},E)\).  That continuity is the special feature of the
concrete \(C(K)\)-predual representation.

\section*{Human Review Recommendation}

The main point to verify is Lemma \ref{lem:factor}: in the case
\(X=C(K)\), the concrete projection \(P\) above is the function-side form of
the antihomogeneous boundary projection \(D\) used by Ondrej--Spurny.  Once
that identification is accepted, the Right-topology argument is formal.

\begin{thebibliography}{9}

\bibitem{ondrej-spurny}
S. Ondrej and J. Spurny,
\emph{Operators on injective tensor products of \(L_1\)-preduals},
arXiv:2604.18157v1, 2026.

\end{thebibliography}

\end{document}
